\section*{Sensitive Metadata Accumulation and Device-Aware Video Anonymization in Endoscopy}

\subsection*{Abstract}
Endoscopic video routinely contains patient-identifying overlays, operator
identifiers, timestamps, and institutional metadata that are clinically useful
during care but unsafe for secondary use without de-identification. This paper
presents a practical anonymization concept for endoscopy environments based on
two linked ideas: (1) sensitive metadata accumulation across sampled frames and
(2) device-aware generation of an anonymized output video. Rather than treating
each frame as an isolated unit, the method aggregates weak and partial evidence
over time into a consolidated metadata representation, then applies a
device-specific anonymization strategy to preserve diagnostically relevant image
content while suppressing protected health information. The approach is designed
for real-world endoscopy infrastructure where overlay geometry, rendering style,
and data density differ by device family.

\subsection*{1. Clinical Motivation}
Modern endoscopy systems superimpose textual and symbolic information directly
onto the live image stream. These overlays often include patient name elements,
date of birth, accession or case identifiers, examination date/time, examiner
identity, and center labels. In routine clinical operation, these fields are
valuable for workflow continuity and medico-legal traceability. However, for
quality improvement, education, AI development, and multicenter research,
retaining such information creates substantial privacy and governance risk.

The central clinical challenge is therefore not only to ``remove text,'' but to
remove identifying context without degrading interpretation of mucosal findings,
instrument position, or procedure dynamics.

\subsection*{2. Conceptual Framework}
\subsubsection*{2.1 Sensitive Metadata Accumulation}
Sensitive metadata accumulation is the process of integrating evidence from many
frames into a unified patient- and procedure-level representation. In endoscopy,
single-frame OCR often produces incomplete fragments due to motion blur,
compression artifacts, transient overlays, and changing luminance. A clinically
useful system must therefore combine partial detections over time, merging
fields when confidence increases and preserving prior high-confidence elements.

This temporal accumulation model improves robustness in two ways:
\begin{itemize}
  \item it reduces dependence on any single frame,
  \item it supports early identification of critical identifiers even when full
  demographic blocks are never fully visible in one image.
\end{itemize}

\subsubsection*{2.2 Device-Aware Anonymized Video Return}
Endoscopy platforms differ in overlay placement, image viewport geometry, and
metadata panel structure. A single fixed anonymization mask is rarely safe
across device types. Device-aware anonymization therefore selects or validates
ROI strategy according to the active endoscopy system. In practical terms, this
means the anonymization layer preserves the endoscope image region and suppresses
side-panel or banner zones where identifying metadata is expected.

The resulting output is not just a ``cleaned'' frame set, but a clinically
usable anonymized video that remains interpretable for diagnosis review and
downstream analytics.

\subsection*{3. Endoscopy Workflow Integration}
In a real endoscopy environment, anonymization must operate under time and
infrastructure constraints:
\begin{itemize}
  \item heterogeneous source videos from different rooms and vendors,
  \item variable codec quality and frame rates,
  \item inconsistent overlay persistence across procedure phases,
  \item requirement for reproducible output suitable for audit.
\end{itemize}

The proposed paradigm aligns with this reality by combining sampled frame
analysis (for efficiency), accumulation-based metadata consolidation (for
clinical reliability), and strategy selection based on device characteristics
(for safer redaction boundaries).

\subsection*{4. Clinical Relevance of Accumulated Metadata}
Accumulated metadata serves three clinical functions beyond redaction itself:
\begin{enumerate}
  \item \textbf{Verification:} confirms whether known identifying fields were
  present in the source material.
  \item \textbf{Auditability:} provides structured evidence for privacy
  compliance and internal governance review.
  \item \textbf{Context Preservation:} enables retention of non-identifying
  procedural descriptors while removing direct identifiers.
\end{enumerate}

This distinction is important: anonymization should minimize patient
re-identification risk while preserving operationally meaningful metadata where
policy allows.

\subsection*{5. Device-Specific Anonymization Strategies}
Two strategy families are clinically relevant in endoscopy:
\begin{itemize}
  \item \textbf{Mask-based viewport preservation:} keep the endoscopic image
  region and suppress surrounding metadata zones according to device geometry.
  \item \textbf{Frame-level suppression or removal:} selectively remove or
  transform frames judged to contain high-risk identifying overlays.
\end{itemize}

Device-aware selection between these strategies can be guided by local policy,
research protocol, and expected downstream use. For teaching archives, viewport
preservation may maximize educational value. For strict secondary-use contexts,
more aggressive suppression may be preferable.

\subsection*{6. Safety, Limitations, and Governance}
Despite strong practical utility, several limitations remain clinically
important:
\begin{itemize}
  \item transient identifiers may be missed under sparse sampling,
  \item overlay drift after firmware updates can invalidate legacy ROI
  assumptions,
  \item OCR uncertainty in low-quality captures can propagate into incomplete
  metadata accumulation.
\end{itemize}

Accordingly, endoscopy anonymization should be treated as a governed clinical
process rather than a one-time technical transform. Recommended safeguards
include routine validation on local device output, periodic policy review with
privacy officers, and release gates based on predefined identifier recall
criteria.

\subsection*{7. Implications for Medical AI and Research}
A reliable accumulation-plus-device-awareness model supports safer scaling of
endoscopy data use in three domains:
\begin{itemize}
  \item multicenter dataset construction with reduced privacy risk,
  \item continuous quality improvement programs using longitudinal video review,
  \item development of AI models on de-identified yet clinically faithful visual
  data.
\end{itemize}

The key contribution is not only improved redaction accuracy, but operational
trust: clinicians and governance teams can understand why a video is considered
anonymized and which metadata signals informed that decision.

\subsection*{8. Results and Evaluation (Illustrative)}
To illustrate practical endpoint evaluation for the accumulation paradigm, a
comparative analysis can be reported against a single-frame baseline using two
clinical priorities: patient-safety protection and preservation of diagnostic
content. A representative reporting template is shown below.

\begin{table}[h]
\centering
\begin{tabular}{|l|c|c|c|}
\hline
\textbf{Metric} & \textbf{Baseline (Single-Frame)} & \textbf{Proposed (Accumulation)} & \textbf{$\Delta$} \\ \hline
PHI Leakage Rate (\%) & 12.4\% & 0.8\% & -93.5\% \\ \hline
Entity Recall (NER) & 0.76 & 0.94 & +23.6\% \\ \hline
Diagnostic Utility (VPR) & 82.1\% & 96.4\% & +17.4\% \\ \hline
\end{tabular}
\caption{Illustrative comparison of anonymization performance. VPR: Viewport Preservation Ratio.}
\end{table}

In this evaluation framing, strict PHI leakage is defined at frame level:
if at least one ground-truth PHI entity is not sufficiently covered by a
predicted anonymization region, that frame is counted as a leakage event. This
strict criterion better reflects clinical privacy expectations than permissive
averaging. Concurrently, VPR quantifies preservation of the diagnostically
relevant endoscopic viewport, preventing over-redaction from being interpreted
as improved safety.

\subsection*{9. Deleting Records Safely}
Safe deletion in endoscopy anonymization should be treated as a controlled
clinical workflow, not a single database action. In practice, this means
combining transaction boundaries, scoped deletion, and auditable operation logs.
The current validation flow demonstrates strong transactional intent and state
tracking, but deletion policy should explicitly preserve those same guarantees.

\subsubsection*{9.1 Operational Principles}
\begin{itemize}
  \item \textbf{Use atomic units of work:} related updates (metadata state,
  annotation cleanup, and parent-record state transitions) should commit or
  rollback together.
  \item \textbf{Delete by scope, not broad criteria:} frame-level annotation
  deletion must be constrained by video, frame range, label, source, and model
  provenance to avoid collateral data loss.
  \item \textbf{Keep a durable audit trail:} every destructive operation should
  emit operation logs with resource type, identifier, actor, and before/after
  status.
  \item \textbf{Prefer idempotent cleanup:} repeated deletion requests should
  converge to the same safe state without introducing inconsistency.
\end{itemize}

\subsubsection*{9.2 Clinical Safety Guardrails}
\begin{itemize}
  \item \textbf{Separation of concerns:} distinguish metadata validation from
  deletion endpoints so clinical review and destructive actions are never
  conflated.
  \item \textbf{Policy gates before hard delete:} sensitive records should pass
  role-based and governance checks before irreversible removal.
  \item \textbf{Retention-aware execution:} deletion logic should respect legal
  retention windows and institutional policy for traceability.
  \item \textbf{Post-delete verification:} safety checks should verify no
  orphaned frame annotations or inconsistent state records remain.
\end{itemize}

Applied in this way, deletion becomes a governed anonymization lifecycle step:
validated metadata and annotation artifacts can be safely removed when policy
requires, without compromising auditability or clinical data integrity.

\subsection*{10. Conclusion}
In endoscopy practice, effective anonymization requires both temporal and
contextual intelligence. Sensitive metadata accumulation across frames provides
temporal robustness, while device-aware anonymized video generation provides
contextual safety aligned with platform-specific overlay behavior. Together,
these principles move de-identification from a narrow technical task toward a
clinically grounded, governance-ready process for secondary use of endoscopic
video data.
